\chapter{Boolean Analysis — Polynomial Approximation}
%%% generated by the blueprint pipeline from TCSlib.BooleanAnalysis.GaussianSpace.Noise.PolynomialApproximation
\section{Overview}

Declarations of the Lean module \texttt{TCSlib.BooleanAnalysis.GaussianSpace.Noise.PolynomialApproximation}.

\section{Declarations}

\begin{definition}[Gaussian polynomial map into $L^p$]
\lean{BooleanAnalysis.GaussianSpace.gaussianPolynomialLp}
\label{BooleanAnalysis.GaussianSpace.gaussianPolynomialLp}
\leanok
\uses{BooleanAnalysis.GaussianSpace.memLp_eval_gaussian}
Let $br^n$ carry standard Gaussian measure, let $bf$ be a field of scalars that is
either $br$ or $bc$, and let $p$ be a finite nonnegative extended real number.
Polynomial evaluation on $rr^n$ defines a $bf$-linear map from
$bf[X_1,
mldots,X_n]$ to $L^p(rr^n,bf)$, sending each polynomial to its equivalence
class as a function under Gaussian measure.
\end{definition}

\begin{lemma}[Exponential of a linear map in the Gaussian polynomial closure]
\lean{BooleanAnalysis.GaussianSpace.exp_linear_mem_gaussianPolynomial_closure}
\label{BooleanAnalysis.GaussianSpace.exp_linear_mem_gaussianPolynomial_closure}
\leanok
\uses{BooleanAnalysis.GaussianSpace.exists_polynomial_eval_linear, BooleanAnalysis.GaussianSpace.gaussianPolynomialLp, BooleanAnalysis.GaussianSpace.memLp_eval_gaussian, BooleanAnalysis.GaussianSpace.tendsto_gaussian_exp_taylor}
\difficulty{5}
Let $n$ be a natural number, let $1\le p<\infty$, let $L\colon\mathbb R^n\to\mathbb C$
be a continuous real-linear map, and let $f$ belong to the complex-valued $L^p$ space
for product standard Gaussian measure on $\mathbb R^n$. If $f(x)=\exp(L(x))$ for
Gaussian-almost every $x$, then $f$ lies in the $L^p$ closure of the set of functions
represented by complex polynomials in the coordinates on $\mathbb R^n$.
\end{lemma}

\begin{lemma}[Gaussian torus Fourier approximation]
\lean{BooleanAnalysis.GaussianSpace.gaussianTorusFourier_mem_polynomial_closure}
\label{BooleanAnalysis.GaussianSpace.gaussianTorusFourier_mem_polynomial_closure}
\leanok
\uses{BooleanAnalysis.GaussianSpace.coeFn_gaussianTorusFourierLp, BooleanAnalysis.GaussianSpace.exp_linear_mem_gaussianPolynomial_closure, BooleanAnalysis.GaussianSpace.gaussianPolynomialLp, BooleanAnalysis.GaussianSpace.gaussianTorusFourierLp, BooleanAnalysis.GaussianSpace.mFourier_gaussianTorusMap}
\difficulty{3}
Let $n$ be a nonnegative integer, let $p$ satisfy $1\le p<\infty$, and let
$m:(\{1,\ldots,n\}\times\{0,1\})\to\mathbb Z$ be an integer frequency vector. The
Fourier character on the product torus with frequency $m$, pulled back along the map
sending $x\in\mathbb R^n$ to the torus coordinates $x_i$ and $\sqrt{2}x_i$ modulo $1$,
belongs in Gaussian $L^p(\mathbb R^n;\mathbb C)$ to the closure of the restrictions of
complex polynomials.
\end{lemma}

\begin{lemma}[Gaussian polynomial lp closure eq top]
\lean{BooleanAnalysis.GaussianSpace.gaussianPolynomialLp_closure_eq_top}
\statementsource{boolean-ch11-gaussian-invariance}{pdf-8b2c50cf9621-p007-b001}
\label{BooleanAnalysis.GaussianSpace.gaussianPolynomialLp_closure_eq_top}
\leanok
\uses{BooleanAnalysis.GaussianSpace.gaussianPolynomialLp, BooleanAnalysis.GaussianSpace.gaussianTorusFourier_mem_polynomial_closure, BooleanAnalysis.GaussianSpace.span_gaussianTorusFourier_dense}
\difficulty{3}
% dataset-pipeline: stub -- prose pending build_informal_statements.py
Statement pending regeneration.
\end{lemma}
